% =============================================================================
% UIDT Framework v3.9 Canonical — arXiv Submission Draft v1
% DOI: 10.5281/zenodo.17835200
% Prepared: 2026-05-24
% Author: P. Rietz
% Status: DRAFT — for external peer-review preparation
% =============================================================================
\documentclass[12pt,a4paper]{article}

% --- Packages ---
\usepackage{amsmath,amssymb,amsthm}
\usepackage{mathtools}
\usepackage{physics}
\usepackage{hyperref}
\usepackage{geometry}
\usepackage{booktabs}
\usepackage{array}
\usepackage{microtype}
\usepackage{xcolor}
\usepackage{cleveref}
\usepackage[numbers,sort&compress]{natbib}

\geometry{margin=2.5cm}

% --- Theorem environments ---
\newtheorem{theorem}{Theorem}[section]
\newtheorem{lemma}[theorem]{Lemma}
\newtheorem{proposition}[theorem]{Proposition}
\newtheorem{corollary}[theorem]{Corollary}
\newtheorem{definition}[theorem]{Definition}
\newtheorem{remark}[theorem]{Remark}

% --- Custom commands ---
\newcommand{\UIDT}{\textsc{UIDT}}
\newcommand{\Lag}{\mathcal{L}}
\newcommand{\kap}{\kappa}
\newcommand{\lams}{\lambda_S}
\newcommand{\gam}{\gamma}
\newcommand{\Dstar}{\Delta^*}
\newcommand{\rvac}{\rho_{\mathrm{vac}}}
\newcommand{\rhobs}{\rho_{\mathrm{vac}}^{\mathrm{obs}}}
\newcommand{\Fmn}{F^a_{\mu\nu}}
\newcommand{\openitem}[1]{\textbf{[OPEN: #1]}}

\hypersetup{
  colorlinks=true,
  linkcolor=blue!60!black,
  citecolor=blue!60!black,
  urlcolor=blue!60!black,
  pdftitle={UIDT Framework v3.9: Yang-Mills Spectral Gap, Scalar Condensate, and Calibrated Cosmological Mapping},
  pdfauthor={P. Rietz}
}

% =============================================================================
\begin{document}

\title{%
  \textbf{Unified Interaction-Density Theory (UIDT) Framework v3.9 Canonical:}\\[6pt]
  Yang-Mills Spectral Gap, Scalar Condensate Coupling,\\[4pt]
  and Falsifiable Cosmological Predictions
}

\author{
  Philipp Rietz\\[4pt]
  \small Independent Researcher\\[2pt]
  \small \href{mailto:p.rietz@uidt-framework.org}{p.rietz@uidt-framework.org}\\[4pt]
  \small Repository: \href{https://github.com/Mass-Gap/UIDT-Framework-v3.9-Canonical}{github.com/Mass-Gap/UIDT-Framework-v3.9-Canonical}\\[2pt]
  \small DOI: \href{https://doi.org/10.5281/zenodo.17835200}{10.5281/zenodo.17835200}
}

\date{May 2026 — Draft v1 (arXiv preparation)}

\maketitle

\begin{abstract}
We present the \UIDT{} Framework v3.9 Canonical, a theoretical construction
coupling a SU(3) Yang--Mills sector to a real scalar field $S$ through
a dimension-five interaction $\frac{\kap}{\Lambda}\,S\,\mathrm{Tr}(FF)\,\Omega$.
The framework yields a Yang--Mills spectral gap
$\Dstar = 1.710 \pm 0.015\,\text{GeV}$, lattice-compatible at $z \approx 0.37\sigma$
(Stratum~B), a calibrated anomalous dimension $\gam = 16.339$ (Stratum~A$^-$)
satisfying the renormalisation-group constraint $5\kap^2 = 3\lams$ with
residual $< 10^{-14}$ (Stratum~A), and a set of falsifiable cosmological
calibrations including $H_0 = 70.4 \pm 0.16\,\text{km\,s}^{-1}\text{Mpc}^{-1}$,
$w_0 = -0.99$, and vacuum energy density
$\rhobs = 2.45 \times 10^{-47}\,\text{GeV}^4$
(all Stratum~C).
Distinct experimental predictions include a Casimir force enhancement
of $+0.59\%$ at $\lambda_{\UIDT} = 0.66\,\text{nm}$ and an X17-floor at
$17.10\,\text{MeV}$ (Stratum~D).
Two mandatory open limitations are explicitly stated:
an unresolved $\sim10^{10}$ residual factor in the vacuum suppression
product (L1) and an unresolved $\delta\gam = 0.0047$ RG correction
(L4), each accompanied by falsification criteria.
All numerical results are reproducible via the public repository
with one command at \texttt{mp.dps\,=\,80} precision.
\end{abstract}

\tableofcontents
\newpage

% =============================================================================
\section{Introduction}
\label{sec:intro}

The cosmological constant problem and the mass-gap conjecture of Yang--Mills
theory represent two of the most significant unresolved problems in theoretical
physics \cite{Weinberg1989,Jaffe2006}.
The \UIDT{} Framework proposes a unified treatment: a scalar field $S$
condensing in the Yang--Mills vacuum generates both a spectral gap and a
suppression mechanism for the QFT vacuum energy density.

The present manuscript documents the \UIDT{} v3.9 Canonical state,
superseding all prior versions (v3.7.1 and earlier).
Version 4.1 adds audit, citation, and numerical-precision layers only;
it does not alter the v3.9 physics.

\paragraph{Scope and limitations.}
This work operates strictly within three evidence strata:
\begin{enumerate}
  \item[I.] Empirical values with stated uncertainties.
  \item[II.] Standard QFT, lattice QCD, and $\Lambda$CDM consensus results.
  \item[III.] \UIDT{} mappings, predictions, and interpretations.
\end{enumerate}
Strata are never silently merged.
Two mandatory open limitations (L1, L4) are carried throughout;
no prestige language is employed.

% =============================================================================
\section{Lagrangian and Parameter Ledger}
\label{sec:lagrangian}

\subsection{Canonical v3.9 Lagrangian}

The \UIDT{} Lagrangian density reads
\begin{equation}
  \Lag_{\UIDT} = -\tfrac{1}{4}\Fmn F^{a\mu\nu}
  + \tfrac{1}{2}(\partial_\mu S)^2 - V(S)
  - \tfrac{\kap}{4}S^2\,\mathrm{Tr}(FF),
  \label{eq:lagrangian_v39}
\end{equation}
where $F^a_{\mu\nu}$ is the SU(3) field-strength tensor and $V(S)$ the
scalar potential.
The v4.1 audit form
\begin{equation}
  \Lag_{\UIDT}^{\mathrm{v4.1}} = -\tfrac{1}{4}\Fmn F^{a\mu\nu}
  + \tfrac{1}{2}\partial_\mu S\partial^\mu S - V(S)
  + \tfrac{\kap}{\Lambda}\,S\,\mathrm{Tr}(FF)\,\Omega
  \label{eq:lagrangian_v41}
\end{equation}
differs dimensionally from Eq.~\eqref{eq:lagrangian_v39}.
This constitutes a \textbf{[TENSION ALERT]}: the two forms are recorded
without silent reconciliation; resolution requires explicit operator matching.

\subsection{Immutable Parameter Ledger}

All parameters below are frozen in \texttt{LEDGER/CLAIMS.json}.

\begin{table}[h!]
  \centering
  \caption{\UIDT{} v3.9 Canonical parameter ledger with evidence tags
           and stratum classification.}
  \label{tab:ledger}
  \renewcommand{\arraystretch}{1.25}
  \begin{tabular}{llll}
    \toprule
    Parameter & Value & Evidence Tag & Stratum \\
    \midrule
    $\Delta / \Dstar$ & $1.710 \pm 0.015\,\text{GeV}$ & [B] & II/III \\
    $\gam$ & $16.339$ & [A$^-$] & III \\
    $\gam_\infty$ & $16.3437$ & [A$^-$] & III \\
    $\delta\gam$ & $0.0047$ & [A$^-$] & III — \openitem{L4} \\
    $\kap$ & $0.500$ & [A] & III \\
    $\lams$ & $5/12$ (exact) & [A] & III \\
    $v$ & $47.7 \pm 0.5\,\text{MeV}$ & [A] & III \\
    $m_S$ & $1.705\,\text{GeV}$ & [B] & II/III \\
    $E_T$ & $2.44\,\text{MeV}$ & [C] & III \\
    $\lambda_{\UIDT}$ & $0.66\,\text{nm}$ & [C] & III \\
    $H_0$ & $70.4 \pm 0.16\,\text{km\,s}^{-1}\text{Mpc}^{-1}$ & [C] & III \\
    $w_0$ & $-0.99$ & [C] & III \\
    $\rhobs$ & $2.45\times10^{-47}\,\text{GeV}^4$ & [C] & III \\
    \bottomrule
  \end{tabular}
\end{table}

% =============================================================================
\section{Renormalisation-Group Constraint}
\label{sec:rg}

The tree-level RG constraint is
\begin{equation}
  5\kap^2 = 3\lams, \qquad
  \left|5\kap^2 - 3\lams\right| < 10^{-14}.
  \label{eq:rg_constraint}
\end{equation}
With $\kap = 1/2$ (exact) and $\lams = 5/12$ (exact):
\begin{equation}
  5 \times (1/2)^2 = 5/4, \qquad
  3 \times (5/12) = 5/4, \qquad
  \text{residual} = 0\ (\text{exact}).
\end{equation}
Numerical verification at \texttt{mp.dps\,=\,80} yields residual
$< 2 \times 10^{-56}$, consistent with exact arithmetic noise at 80-digit
precision. This constitutes Evidence [A].

\paragraph{Open limitation L4.}
The calibrated value $\gam = 16.339$ carries correction
$\delta\gam = \gam - 49/3 = 17/3000 \approx 0.0056\overline{6}$
(ledger: $0.0047 \pm 0.0015$).
All minimal local quantum-field-theory paths generating this correction
have been systematically excluded in the Phase-8 P1 programme
(PRs \#459--\#498 in the repository).
The BMW/Dyson/FRG non-perturbative operator-mixing route remains the
sole admissible path and is currently open.
\openitem{L4}

% =============================================================================
\section{Yang--Mills Spectral Gap}
\label{sec:gap}

The spectral gap $\Dstar = 1.710 \pm 0.015\,\text{GeV}$ is identified with
the lowest SU(3) glueball mass in the scalar channel.
This is a \UIDT{} Stratum~III mapping; the parameter is \emph{not} a direct
particle observation.

\paragraph{Lattice compatibility.}
The value is compatible with lattice QCD glueball-mass estimates
$m_G \approx 1.5$--$1.8\,\text{GeV}$ at the $z \approx 0.37\sigma$ level [B].
The X2370 candidate state is listed as a Stratum~D prediction.

\paragraph{Banach contraction bound.}
The \UIDT{} construction requires the Banach contraction operator norm
$\|\mathcal{B}\| < 1$. Violation of this bound constitutes a falsification
trigger. Current numerical verification yields $\|\mathcal{B}\| = L$
with $L < 0.9$ (exact value in \texttt{verification/scripts/}).

% =============================================================================
\section{Vacuum Energy Suppression}
\label{sec:vacuum}

The \UIDT{} vacuum energy density satisfies
\begin{equation}
  \rvac^{\mathrm{obs}} =
  \rvac^{\mathrm{QFT}} \times \pi^{-2}
  \times \prod_{n=1}^{99} f_n(g),
  \label{eq:vacuum}
\end{equation}
where $f_n(g)$ are spectral suppression factors derived from the
\UIDT{} coupling $g$.

\paragraph{Open limitation L1.}
The functional form of $f_n(g)$ is not yet derived from first principles.
The current best estimate, using $f_n = 1$ (placeholder), yields a
residual factor $\sim 10^{10}$ relative to $\rhobs$.
This factor is carried as a mandatory open limitation and does not
constitute a resolution of the cosmological constant problem.
The mapping is classified Stratum~C (calibrated cosmology).
\openitem{L1}

% =============================================================================
\section{Torsion Sector}
\label{sec:torsion}

The torsion constraint requires
\begin{equation}
  E_T = 0 \implies \Sigma_T = 0\quad (\text{exact}).
  \label{eq:torsion}
\end{equation}
With $E_T = 2.44\,\text{MeV}$ [C], $\Sigma_T \neq 0$ is admissible.
If an independent measurement establishes $E_T = 0$ while $\Sigma_T \neq 0$
remains, this triggers \texttt{[TORSION\_CONSTRAINT\_FAIL]} and constitutes
a kill-switch.

% =============================================================================
\section{Experimental Predictions}
\label{sec:predictions}

All predictions carry Stratum~D (falsifiable, pending independent verification).

\begin{table}[h!]
  \centering
  \caption{\UIDT{} v3.9 falsifiable predictions.}
  \label{tab:predictions}
  \renewcommand{\arraystretch}{1.3}
  \begin{tabular}{p{4cm}p{3.5cm}p{2.5cm}p{3.5cm}}
    \toprule
    Prediction & Value & Tag & Kill-Switch Condition \\
    \midrule
    Casimir force enhancement at $\lambda = 0.66\,\text{nm}$
      & $+0.59\%$ & [D] & $|\Delta F/F| < 0.1\%$ \\
    X17 floor energy
      & $17.10\,\text{MeV}$ & [D] & Confirmed absence of state \\
    X2370 glueball candidate
      & Stratum D mapping & [D] & Lattice excludes $\Dstar$ by $> 3\sigma$ \\
    Photonic refractive index anomaly
      & $n = 16.339$ & [D] & Null result at required precision \\
    \bottomrule
  \end{tabular}
\end{table}

% =============================================================================
\section{Mandatory Open Limitations}
\label{sec:limitations}

Four limitations are mandatory disclosures and must be carried in any
publication of this framework.

\begin{description}
  \item[L1 — Vacuum suppression.]
    An unresolved factor $\sim 10^{10}$ remains between the \UIDT{} vacuum
    suppression product and $\rhobs$. The f_n(g) functional form is not derived.
    Falsification: Casimir measurement at $\lambda = 0.66\,\text{nm}$ with
    $|\Delta F/F| < 0.1\%$.
  \item[L2 — Electron mass.]
    A $\approx 23\%$ discrepancy exists between the \UIDT{} electron-mass
    prediction and the measured value. Origin not identified.
  \item[L4 — $\delta\gam$ RG-gap.]
    The correction $\delta\gam = 0.0047$ to $\gam = 16.339$ is not derived
    from first principles. All minimal local mechanisms excluded (Phase-8 P1).
    Falsification: Photonic null result at $n = 16.339$.
  \item[L-$\beta$ — Physical $\beta$-functions.]
    Physical one-loop $\beta$-functions are not derived from $\Lag_{\UIDT}$.
    The tree-level RG constraint \eqref{eq:rg_constraint} is exact but does not
    constitute loop-level RG flow.
\end{description}

% =============================================================================
\section{Reproducibility}
\label{sec:repro}

All results in this paper are reproducible via:
\begin{verbatim}
git clone https://github.com/Mass-Gap/UIDT-Framework-v3.9-Canonical
cd UIDT-Framework-v3.9-Canonical
python verification/scripts/verify_all.py
\end{verbatim}
Expected output: all residuals $< 10^{-14}$. Platform: Python 3, \texttt{mpmath}
with \texttt{mp.dps\,=\,80} localised per function. No mock objects.

% =============================================================================
\section{Conclusion}
\label{sec:conclusion}

The \UIDT{} Framework v3.9 Canonical provides a mathematically consistent
construction coupling a Yang--Mills spectral gap to a scalar condensate.
The tree-level RG constraint is exact [A]; the spectral gap is lattice-compatible
at the $0.37\sigma$ level [B]; cosmological calibrations are internally consistent
[C]; and four falsifiable experimental predictions are registered [D].

Two mandatory open limitations (L1 and L4) prevent promotion of the framework
beyond its current evidence level. These are explicitly documented with
falsification criteria. Resolution of L4 via BMW/Dyson/FRG computation and
resolution of L1 via derivation of $f_n(g)$ constitute the primary theoretical
objectives for the next development phase.

% =============================================================================
\appendix

\section{DOI / arXiv Resolvability Table}
\label{app:dois}

\begin{table}[h!]
  \centering
  \renewcommand{\arraystretch}{1.25}
  \begin{tabular}{lll}
    \toprule
    DOI / arXiv & Status & Used for \\
    \midrule
    \href{https://doi.org/10.5281/zenodo.17835200}{10.5281/zenodo.17835200}
      & ✓ Resolved & Primary project DOI \\
    \href{https://zenodo.org/records/18072470}{zenodo.org/18072470}
      & ✓ Resolved & Supplementary results \\
    \href{https://zenodo.org/records/18740600}{zenodo.org/18740600}
      & ✓ Resolved & Claims addendum \\
    \href{https://zenodo.org/records/19228489}{zenodo.org/19228489}
      & ✓ Resolved & Verification logs \\
    \bottomrule
  \end{tabular}
  \caption{All Zenodo DOIs verified resolvable as of 2026-05-24.}
  \label{tab:dois}
\end{table}

% --- Bibliography ---
\begin{thebibliography}{99}

\bibitem{Weinberg1989}
S.~Weinberg,
\textit{The cosmological constant problem},
Rev.\ Mod.\ Phys.\ \textbf{61}, 1 (1989).
\href{https://doi.org/10.1103/RevModPhys.61.1}{doi:10.1103/RevModPhys.61.1}

\bibitem{Jaffe2006}
A.~Jaffe and E.~Witten,
\textit{Quantum Yang--Mills theory},
in: The Millennium Prize Problems,
CMI/AMS (2006), pp.\ 129--152.
\href{https://www.claymath.org/millennium-problems}{claymath.org}

\bibitem{Morningstar1999}
C.~J.~Morningstar and M.~Peardon,
\textit{Glueball spectrum from an anisotropic lattice study},
Phys.\ Rev.\ D \textbf{60}, 034509 (1999).
\href{https://doi.org/10.1103/PhysRevD.60.034509}{doi:10.1103/PhysRevD.60.034509}

\bibitem{Pawlowski2021}
J.~M.~Pawlowski et al.,
\textit{Nonperturbative renormalization group approaches},
Phys.\ Rep.\ \textbf{910}, 1 (2021).
\href{https://doi.org/10.1016/j.physrep.2021.01.001}{doi:10.1016/j.physrep.2021.01.001}

\bibitem{Litim2001}
D.~F.~Litim,
\textit{Optimised renormalisation group flows},
Phys.\ Rev.\ D \textbf{64}, 105007 (2001).
\href{https://arxiv.org/abs/hep-th/0103195}{arXiv:hep-th/0103195}

\end{thebibliography}

\end{document}
% =============================================================================
% END OF DRAFT — version history tracked in git
% =============================================================================
